% ============================================================
%  Introducing Python, chapter 22: Data in Space: Networks
% ============================================================
\chapter{Data in Space: Networks}

\tikzset{nb/.style={draw=cblue,fill=cblue!10,rounded corners=2pt,font=\small\sffamily,
  inner sep=4pt,align=center,minimum height=6.5mm}}

Chapter 21 was about doing many things at once. This chapter is about doing things in
many places: programs on different machines that exchange bytes. You get the layers
under every web and LLM API call (TCP/IP, sockets), the common message patterns, and the
tools that sit on top (serialization, RPC, cloud basics).

% ------------------------------------------------------------
\section{TCP/IP}

Networks work because everyone follows layered rules called \emph{protocols}. A layer
only needs to know the layer directly below it.

%MD
%```mermaid
%flowchart TB
%    A["Application<br/>HTTP, SMTP, DNS: your API call lives here"]
%    T["Transport<br/>TCP or UDP: ports, bytes in order (TCP)"]
%    I["Internet<br/>IP: addresses, packets hop between routers"]
%    L["Link<br/>Ethernet, Wi-Fi: electrical and radio signals"]
%    A -->|uses| T -->|uses| I -->|uses| L
%```
%ENDMD
\begin{center}
\begin{tikzpicture}[lay/.style={draw,rounded corners=2pt,minimum width=4.2cm,
    minimum height=8mm,font=\small\sffamily\bfseries,align=center}]
  \node[lay,draw=cgreen,fill=cgreen!12]  (a) at (0,3)  {Application};
  \node[lay,draw=cblue,fill=cblue!12]    (t) at (0,2)  {Transport};
  \node[lay,draw=cpurple,fill=cpurple!12](i) at (0,1)  {Internet};
  \node[lay,draw=cgray,fill=cgray!12]    (l) at (0,0)  {Link};
  \foreach \n/\txt in {a/{HTTP, SMTP, DNS: your API call lives here},
                       t/{TCP or UDP: ports, bytes in order (TCP)},
                       i/{IP: addresses, packets hop between routers},
                       l/{Ethernet, Wi-Fi: electrical and radio signals}}
    \node[anchor=west,font=\footnotesize\sffamily] at ($(\n.east)+(4mm,0)$) {\txt};
  \draw[->,thick,cgray] (-3.2,3.4) -- (-3.2,-0.4);
  \node[rotate=90,font=\scriptsize\sffamily,text=cgray] at (-3.5,1.5) {each layer uses the one below};
\end{tikzpicture}
\end{center}

Two protocols matter at the transport layer.

\begin{center}
\begin{tabularx}{\linewidth}{@{}l X X@{}}
\toprule
 & \textbf{UDP} & \textbf{TCP} \\
\midrule
Unit of data      & datagram (one small packet)       & stream of bytes \\
Setup             & none, just send                   & handshake first (connect) \\
Delivery          & may be lost, duplicated, reordered & retried, in order, no duplicates \\
Speed             & fast, light                       & slower \\
Typical use       & voice, games, DNS lookups         & web, databases, most APIs \\
\bottomrule
\end{tabularx}
\end{center}

HTTP/3 runs on QUIC, which itself sits on UDP, so the choice is not always black and white.

Your own machine is always \code{127.0.0.1}, also called \code{localhost} or the
\emph{loopback} interface. Two programs on one machine can talk over it without any
real network. Every example below uses it.

% ------------------------------------------------------------
\section{Sockets}

A \emph{socket} is the lowest-level network tool: an endpoint you send bytes into and
read bytes from. It comes from C and Unix. Its address is a tuple
\code{(host, port)}: the host is like a building, the port like a door number.

\subsection{UDP: send one datagram}

This example runs server and client in one script. The server runs in a thread (Chapter
21). Port \code{0} asks the OS to pick any free port, so the example never clashes with
another program.

\begin{pycode}
import socket
import threading

ready = threading.Event()
addr = {}

def server():
    s = socket.socket(socket.AF_INET, socket.SOCK_DGRAM)   # DGRAM = UDP
    s.bind(("127.0.0.1", 0))          # port 0: OS picks a free port
    addr["value"] = s.getsockname()
    ready.set()
    data, client = s.recvfrom(4096)   # waits for one datagram
    print("server got:", data)
    s.sendto(data.upper(), client)    # reply to whoever sent it
    s.close()

t = threading.Thread(target=server)
t.start()
ready.wait()

c = socket.socket(socket.AF_INET, socket.SOCK_DGRAM)
c.sendto("ping from eval worker".encode("utf-8"), addr["value"])
reply, who = c.recvfrom(4096)
print("client got:", reply)
print(type(reply), who == addr["value"])
c.close()
t.join()
\end{pycode}
\begin{rawcode}
server got: b'ping from eval worker'
client got: b'PING FROM EVAL WORKER'
<class 'bytes'> True
\end{rawcode}

\begin{typecheck}
Sockets move \code{bytes}, never \code{str}. \code{sendto("text", addr)} raises
\code{TypeError}. Encode on the way out (\code{.encode("utf-8")}) and decode on the way
in (\code{.decode("utf-8")}). \code{AF\_INET} means an IP socket; \code{SOCK\_DGRAM} means
UDP.
\end{typecheck}

The client never calls \code{bind}. The OS gives it a random port, and the server learns
that port from \code{recvfrom}, which returns \code{(bytes, (host, port))}.

\subsection{TCP: a connection}

TCP adds a handshake, so the roles are more visible: the server \emph{listens} and
\emph{accepts}; the client \emph{connects}.

%MD
%```mermaid
%flowchart TB
%    subgraph Server
%        S1["socket()"] --> S2["bind(addr)"] --> S3["listen()"] --> S4["accept() waits"] --> S5["recv() / sendall()"] --> S6["close()"]
%    end
%    subgraph Client
%        C1["socket()"] --> C2["connect(addr)"] --> C3["sendall() / recv()"] --> C4["close()"]
%    end
%    C2 <-->|handshake| S4
%    C3 <-->|bytes| S5
%```
%ENDMD
\begin{center}
\begin{tikzpicture}[st/.style={nb,minimum width=3.4cm},>=Stealth]
  \node[font=\sffamily\bfseries,text=cgreen] at (0,0.7) {Server};
  \node[font=\sffamily\bfseries,text=cblue]  at (7,0.7) {Client};
  \node[st] (s1) at (0,0)  {\code{socket()}};
  \node[st] (s2) at (0,-1) {\code{bind(addr)}};
  \node[st] (s3) at (0,-2) {\code{listen()}};
  \node[st,fill=cgold!15,draw=cgold] (s4) at (0,-3) {\code{accept()} waits};
  \node[st] (s5) at (0,-4) {\code{recv()} / \code{sendall()}};
  \node[st] (s6) at (0,-5) {\code{close()}};
  \node[st] (c1) at (7,-2) {\code{socket()}};
  \node[st,fill=cgold!15,draw=cgold] (c2) at (7,-3) {\code{connect(addr)}};
  \node[st] (c3) at (7,-4) {\code{sendall()} / \code{recv()}};
  \node[st] (c4) at (7,-5) {\code{close()}};
  \foreach \a/\b in {s1/s2,s2/s3,s3/s4,s4/s5,s5/s6,c1/c2,c2/c3,c3/c4}
    \draw[->,cgray] (\a) -- (\b);
  \draw[<->,thick,cgold] (c2.west) -- node[above,font=\scriptsize\sffamily]{handshake} (s4.east);
  \draw[<->,thick,cteal] (c3.west) -- node[above,font=\scriptsize\sffamily]{bytes} (s5.east);
\end{tikzpicture}
\end{center}

TCP code looks easy until you meet one problem: where does one message end?

\subsection{TCP is a stream, not messages}

\begin{boss}[: no message boundaries]
TCP delivers a \emph{stream of bytes}. If you call \code{sendall} three times, the other
side may get all three pieces in one \code{recv}, or one piece split into two. The
number of \code{send} calls and the number of \code{recv} calls are unrelated.
\end{boss}

%MD
%```text
%byte stream:  H E L L O N E T W O R L D
%sends      :  [ send 1 ][ s2 ][ send 3 ]      (5 + 3 + 5 bytes)
%recvs      :  [ recv 1: 8 B  ][r2:3][r3]      (8 + 3 + 2 bytes)
%```
%ENDMD
\begin{center}
\begin{tikzpicture}[cell/.style={draw=cgray,fill=white,minimum width=5.6mm,minimum height=6mm,
    inner sep=0,font=\ttfamily\footnotesize},>=Stealth]
  \foreach \i/\c in {1/H,2/E,3/L,4/L,5/O,6/N,7/E,8/T,9/W,10/O,11/R,12/L,13/D}
    \node[cell] at (\i*0.56,0) {\c};
  \foreach \a/\b/\t in {1/5/send 1,6/8/send 2,9/13/send 3}
    \draw[decorate,decoration={brace,amplitude=3pt},cblue,thick]
      ({\a*0.56-0.28},0.35) -- node[above=2pt,font=\scriptsize\sffamily,text=cblue]{\t}
      ({\b*0.56+0.28},0.35);
  \foreach \a/\b/\t in {1/8/recv 1: 8 bytes,9/11/recv 2: 3,12/13/recv 3: 2}
    \draw[decorate,decoration={brace,amplitude=3pt,mirror},cred,thick]
      ({\a*0.56-0.28},-0.35) -- node[below=2pt,font=\scriptsize\sffamily,text=cred]{\t}
      ({\b*0.56+0.28},-0.35);
\end{tikzpicture}
\end{center}

To exchange whole messages you must add structure yourself: a fixed size, a delimiter
character, or a length prefix. The length prefix is the most robust. Write a 4-byte
number first, then the body.

%MD
%```text
%+-------------+-----------------+
%| 00 00 00 05 | H E L L O       |
%+-------------+-----------------+
%  4-byte length   body: exactly 5 bytes
%  (big-endian)
%```
%ENDMD
\begin{center}
\begin{tikzpicture}
  \node[draw=cgold,fill=cgold!15,minimum width=2.6cm,minimum height=7mm,font=\ttfamily\small]
    (h) at (0,0) {00 00 00 05};
  \node[draw=cgreen,fill=cgreen!12,minimum width=3.6cm,minimum height=7mm,font=\ttfamily\small,
    right=0pt of h] (b) {HELLO};
  \node[font=\scriptsize\sffamily,text=cgold,below=1mm of h] {4-byte length (big-endian)};
  \node[font=\scriptsize\sffamily,text=cgreen,below=1mm of b] {body: exactly 5 bytes};
\end{tikzpicture}
\end{center}

\begin{pycode}
import socket
import struct
import threading

ready = threading.Event()
box = {}

def recv_exact(sock, n):
    """Read exactly n bytes, however many recv() calls it takes."""
    chunks = []
    while n > 0:
        chunk = sock.recv(n)
        if not chunk:                      # b"" means the peer closed
            raise ConnectionError("peer closed early")
        chunks.append(chunk)
        n -= len(chunk)
    return b"".join(chunks)

def send_msg(sock, text):
    body = text.encode("utf-8")
    sock.sendall(struct.pack(">I", len(body)) + body)  # header + body

def recv_msg(sock):
    (size,) = struct.unpack(">I", recv_exact(sock, 4))
    return recv_exact(sock, size).decode("utf-8")

def server():
    with socket.socket(socket.AF_INET, socket.SOCK_STREAM) as srv:
        srv.bind(("127.0.0.1", 0))
        srv.listen()
        box["addr"] = srv.getsockname()
        ready.set()
        conn, peer = srv.accept()          # blocks until a client connects
        with conn:
            while True:
                try:
                    msg = recv_msg(conn)
                except ConnectionError:
                    break
                send_msg(conn, "echo: " + msg)

t = threading.Thread(target=server)
t.start()
ready.wait()

with socket.socket(socket.AF_INET, socket.SOCK_STREAM) as cli:
    cli.connect(box["addr"])
    for prompt in ["hello", "second request", "x" * 5000]:
        send_msg(cli, prompt)
        reply = recv_msg(cli)
        print(len(prompt), "->", reply[:22], len(reply))
t.join()
print(struct.pack(">I", 5))
\end{pycode}
\begin{rawcode}
5 -> echo: hello 11
14 -> echo: second request 20
5000 -> echo: xxxxxxxxxxxxxxxx 5006
b'\x00\x00\x00\x05'
\end{rawcode}

The 5000-byte message needs several \code{recv} calls underneath. Because
\code{recv\_exact} loops, the caller never notices.

\begin{trap}[: \code{recv(n)} may return fewer than \(n\) bytes]
\code{recv(1000)} means \emph{at most} 1000. Assuming you got a whole message is the
classic socket bug: it works on localhost with tiny messages and fails in production.
An empty result \code{b""} is not ``no data yet''; it means the other side closed.
\end{trap}

\begin{typecheck}
Length is counted in \emph{bytes}, not characters. \code{len("caf\textbackslash u00e9")} is
4 but \code{len("caf\textbackslash u00e9".encode("utf-8"))} is 5. Compute the header from the
encoded bytes, as \code{send\_msg} does. \code{struct.pack(">I", 5)} gives the \code{bytes}
value \code{b'\textbackslash x00\textbackslash x00\textbackslash x00\textbackslash x05'}.
\end{typecheck}

\begin{trap}[: ``Address already in use'']
A fixed port such as 6789 stays busy while an old server (or a recently closed
connection) holds it. Use port 0 for tests, or kill the old process.
\end{trap}

\begin{gamerule}[: bind to localhost while learning]
Use \code{"127.0.0.1"} in \code{bind}. Binding to \code{"0.0.0.0"} or \code{"*"} exposes
the port to every network the machine is on. Do that only on purpose.
\end{gamerule}

\subsection{Scapy and Netcat}

\begin{center}
\begin{tabularx}{\linewidth}{@{}l X@{}}
\toprule
\textbf{Tool} & \textbf{What it is for} \\
\midrule
Scapy (\code{pip install scapy}) & build, send and inspect raw packets from Python. Like a
programmable Wireshark. Do not confuse it with Scrapy (web scraping). \\
Netcat (\code{nc}) & command-line tool to open a connection to a host and port and type
bytes by hand. Handy to check that a port answers. \\
\bottomrule
\end{tabularx}
\end{center}

% ------------------------------------------------------------
\section{Networking Patterns}

Real systems are built from a few repeated shapes.

%MD
%```mermaid
%flowchart LR
%    subgraph RR["Request-reply"]
%        c1["client"] -->|request| s1["server"]
%        s1 -->|reply| c1
%    end
%    subgraph PUSH["Push (fan-out)"]
%        p["producer"] -->|job A| w1["worker 1"]
%        p -->|job B| w2["worker 2"]
%    end
%    subgraph PULL["Pull (fan-in)"]
%        a["proc 1"] --> l["logger"]
%        b["proc 2"] --> l
%    end
%    subgraph PS["Publish-subscribe"]
%        pub["publisher"] -->|copy| u1["sub (topic X)"]
%        pub -->|copy| u2["sub (topic X)"]
%    end
%```
%ENDMD
\begin{center}
\begin{tikzpicture}[n/.style={nb,minimum width=1.5cm,font=\scriptsize\sffamily},>=Stealth]
  % request-reply
  \node[font=\footnotesize\sffamily\bfseries] at (1.6,1.6) {Request-reply};
  \node[n] (c) at (0,0.5) {client};
  \node[n,draw=cgreen,fill=cgreen!10] (s) at (3.2,0.5) {server};
  \draw[->,cblue] ([yshift=2pt]c.east) -- ([yshift=2pt]s.west);
  \draw[<-,cgreen] ([yshift=-2pt]c.east) -- ([yshift=-2pt]s.west);
  % push-pull
  \node[font=\footnotesize\sffamily\bfseries] at (8.4,1.6) {Push (fan-out)};
  \node[n] (p) at (6.6,0.5) {producer};
  \node[n,draw=cgreen,fill=cgreen!10] (w1) at (9.8,1.0) {worker 1};
  \node[n,draw=cgreen,fill=cgreen!10] (w2) at (9.8,0.0) {worker 2};
  \draw[->] (p) -- node[above,font=\tiny]{job A} (w1);
  \draw[->] (p) -- node[below,font=\tiny]{job B} (w2);
  % pull
  \node[font=\footnotesize\sffamily\bfseries] at (1.6,-1.2) {Pull (fan-in)};
  \node[n] (a) at (0,-2.2) {proc 1};
  \node[n] (b) at (0,-3.0) {proc 2};
  \node[n,draw=cgreen,fill=cgreen!10] (l) at (3.2,-2.6) {logger};
  \draw[->] (a) -- (l);
  \draw[->] (b) -- (l);
  % pub-sub
  \node[font=\footnotesize\sffamily\bfseries] at (8.4,-1.2) {Publish-subscribe};
  \node[n,draw=cgold,fill=cgold!15] (pub) at (6.6,-2.6) {publisher};
  \node[n,draw=cgreen,fill=cgreen!10] (u1) at (9.8,-2.0) {sub (topic X)};
  \node[n,draw=cgreen,fill=cgreen!10] (u2) at (9.8,-3.0) {sub (topic X)};
  \draw[->] (pub) -- node[above,font=\tiny]{copy} (u1);
  \draw[->] (pub) -- node[below,font=\tiny]{copy} (u2);
\end{tikzpicture}
\end{center}

\begin{center}
\begin{tabularx}{\linewidth}{@{}l X X@{}}
\toprule
\textbf{Pattern} & \textbf{Who gets a message} & \textbf{Example} \\
\midrule
Request-reply & the one server asked; client waits & browser to web server, call to an LLM API \\
Push (fan-out) & exactly one worker from a pool & job queue, load balancer \\
Pull (fan-in) & one collector reads many sources & log collector \\
Pub-sub & \emph{every} subscriber of the topic; none means dropped & event notifications \\
\bottomrule
\end{tabularx}
\end{center}

Push and pub-sub look similar but differ in one question: does the message go to one
receiver or to all of them? Both shapes in plain Python, with no network:

\begin{pycode}
import queue
import threading

# push / pull: one queue, several workers; each job goes to ONE worker
jobs = queue.Queue()
done = queue.Queue()

def worker(name):
    while True:
        job = jobs.get()
        if job is None:
            break
        done.put((name, job, job * job))

pool = [threading.Thread(target=worker, args=(f"w{i}",)) for i in range(2)]
for t in pool:
    t.start()
for n in range(4):
    jobs.put(n)
for _ in pool:
    jobs.put(None)               # one stop signal per worker
for t in pool:
    t.join()
results = sorted(done.get()[1:] for _ in range(4))
print(results)

# pub-sub: every subscriber of a topic gets its OWN copy
class Broker:
    def __init__(self):
        self.subs = {}           # topic -> list of callbacks

    def subscribe(self, topic, callback):
        self.subs.setdefault(topic, []).append(callback)

    def send(self, topic, message):
        for callback in self.subs.get(topic, []):   # no subscriber: dropped
            callback(message)

bus = Broker()
bus.subscribe("eval.done", lambda m: print("logger  :", m))
bus.subscribe("eval.done", lambda m: print("alerter :", m))
bus.subscribe("eval.failed", lambda m: print("pager   :", m))
bus.send("eval.done", {"run": 7, "score": 0.91})
bus.send("eval.cost", {"usd": 1.2})        # nobody listens: vanishes
\end{pycode}
\begin{rawcode}
[(0, 0), (1, 1), (2, 4), (3, 9)]
logger  : {'run': 7, 'score': 0.91}
alerter : {'run': 7, 'score': 0.91}
\end{rawcode}

% ------------------------------------------------------------
\section{The Request-Reply Pattern}

Request-reply is the everyday pattern: you ask DNS, a web server or an LLM API, and you
get an answer or an error.

\subsection{ZeroMQ}

ZeroMQ (package \code{pyzmq}) is a library, not a server. It wraps sockets and gives you
what plain sockets do not: whole messages (no framing code), automatic reconnects, and
queued messages when the peer is late. The ``zero'' means no broker is needed. You can
start client and server in either order.

\begin{center}
\begin{tabularx}{\linewidth}{@{}l l X@{}}
\toprule
\textbf{Socket type} & \textbf{Pattern} & \textbf{Note} \\
\midrule
\code{REQ} / \code{REP} & synchronous request / reply & strict send, receive, send, receive \\
\code{DEALER} / \code{ROUTER} & asynchronous request / reply & build brokers between many clients and servers \\
\code{PUSH} / \code{PULL} & fan-out / fan-in & ``ventilator'': farm out tasks, gather results \\
\code{PUB} / \code{SUB} & publish / subscribe & subscribers filter by topic prefix \\
\bottomrule
\end{tabularx}
\end{center}

Needs \code{pip install pyzmq}. I could not run these here (no package installed); the
files were syntax-checked only.

\begin{pycode}
# server.py   (needs: pip install pyzmq)
import zmq

ctx = zmq.Context()
server = ctx.socket(zmq.REP)
server.bind("tcp://127.0.0.1:5555")
while True:
    request = server.recv()                    # bytes, one whole message
    text = request.decode("utf-8")
    server.send(("tokens: %d" % len(text.split())).encode("utf-8"))
\end{pycode}
\begin{pycode}
# client.py
import zmq

ctx = zmq.Context()
client = ctx.socket(zmq.REQ)
client.connect("tcp://127.0.0.1:5555")         # connect, not bind
client.send("count my words please".encode("utf-8"))
print(client.recv().decode("utf-8"))           # tokens: 4
\end{pycode}

The server \emph{binds} (it is the stable side); the client \emph{connects}. Many REQ
clients can connect to one REP server; it handles them one at a time and ZeroMQ queues
the rest. Strings must be encoded: ZeroMQ messages are bytes. For structured data, send
JSON bytes.

ZeroMQ scales by changing only the address scheme:

\begin{center}
\begin{tabular}{@{}l l@{}}
\toprule
\textbf{Transport} & \textbf{Connects} \\
\midrule
\code{tcp://host:port} & processes, same or different machines \\
\code{ipc://name} & processes on one machine \\
\code{inproc://name} & threads inside one process (no locks needed) \\
\bottomrule
\end{tabular}
\end{center}

\subsection{Other messaging tools}

\begin{center}
\begin{tabularx}{\linewidth}{@{}l X@{}}
\toprule
\textbf{Tool} & \textbf{Notes} \\
\midrule
RabbitMQ (Python client: Pika) & popular message broker with queues and routing \\
ActiveMQ (STOMP protocol) & Apache broker, several Python clients \\
NATS & fast, light messaging; request-reply, pub-sub, queues \\
\bottomrule
\end{tabularx}
\end{center}

% ------------------------------------------------------------
\section{The Publish-Subscribe Pattern}

Pub-sub is a broadcast, not a queue. Publishers announce messages on a \emph{topic}; each
subscriber says which topics it wants. A message can be processed once, many times or
never. A publisher does not know who listens.

%MD
%```mermaid
%flowchart LR
%    P["publisher"] --> T["topics: eval.done, eval.cost"]
%    T --> A["logger: eval.done"]
%    T --> B["alerter: eval.done"]
%    T -.-> C["nobody: eval.cost"]
%```
%ENDMD
\begin{center}
\begin{tikzpicture}[n/.style={nb,minimum width=2.2cm,font=\scriptsize\sffamily},>=Stealth]
  \node[n,draw=cgold,fill=cgold!15] (p) at (0,0) {publisher};
  \node[n,draw=cpurple,fill=cpurple!10,minimum width=2.6cm] (t) at (3.6,0)
    {topics: \code{eval.done}, \code{eval.cost}};
  \node[n,draw=cgreen,fill=cgreen!10] (a) at (8,0.9) {logger: \code{eval.done}};
  \node[n,draw=cgreen,fill=cgreen!10] (b) at (8,0) {alerter: \code{eval.done}};
  \node[n,draw=cgray,fill=cgray!10] (c) at (8,-0.9) {nobody: \code{eval.cost}};
  \draw[->] (p) -- (t);
  \draw[->] (t) -- (a);
  \draw[->] (t) -- (b);
  \draw[->,dashed,cgray] (t) -- (c);
\end{tikzpicture}
\end{center}

\subsection{Redis}

Redis (Chapter 23 covers it as a data store) also has pub-sub. Start the subscriber
\emph{first}: messages sent while nobody is subscribed are lost. Needs
\code{pip install redis} and a running Redis or Valkey server; syntax-checked only.

\begin{pycode}
# needs: pip install redis, and a Redis (or Valkey) server on localhost
import redis

r = redis.Redis(decode_responses=True)         # return str, not bytes

sub = r.pubsub()
sub.subscribe("eval.done")
for msg in sub.listen():
    if msg["type"] == "message":
        print(msg["channel"], msg["data"])
        break

# in another process:  r.publish("eval.done", "run 7")

# push / pull with a blocking pop: waits up to 2 s for a job
r.lpush("jobs", "prompt-1", "prompt-2")
item = r.brpop("jobs", timeout=2)              # (list_name, value) or None
print(item)
\end{pycode}

\begin{typecheck}
Without \code{decode\_responses=True}, redis-py returns \code{bytes} (\code{b'run 7'}),
so comparisons with \code{str} silently fail. \code{brpop} returns a tuple
\code{(key, value)}, or \code{None} on timeout.
\end{typecheck}

\subsection{ZeroMQ}

ZeroMQ pub-sub has no central server: each publisher writes to its connected
subscribers. Messages are \emph{multipart}: the first part is the topic.

\begin{pycode}
# needs: pip install pyzmq
import time
import zmq

ctx = zmq.Context()
pub = ctx.socket(zmq.PUB)
pub.bind("tcp://127.0.0.1:5556")
time.sleep(1)                                  # let subscribers connect
pub.send_multipart([b"eval.done", b'{"run": 7}'])

sub = ctx.socket(zmq.SUB)
sub.connect("tcp://127.0.0.1:5556")
sub.setsockopt(zmq.SUBSCRIBE, b"eval.")        # topic prefix, as bytes
topic, body = sub.recv_multipart()
print(topic.decode("utf-8"), body.decode("utf-8"))
\end{pycode}

\begin{trap}[: a SUB socket receives nothing by default]
You must subscribe. To get everything, subscribe to the empty bytes \code{b""}. Also,
a PUB that sends right after \code{bind} may lose the first messages because subscribers
have not connected yet (the \code{sleep} above).
\end{trap}

Other options: RabbitMQ with Pika, NATS, PyPubSub (in-process), PubSubHubbub (webhook style).

% ------------------------------------------------------------
\section{Internet Services}

\subsection{Domain Name System (DNS)}

DNS turns names into IP addresses. When a request hangs on ``looking up host'', DNS is
the first suspect.

%MD
%```mermaid
%flowchart LR
%    A["your code<br/>getaddrinfo('api.x.test')"] -->|name?| B["DNS servers<br/>(distributed database)"]
%    B -->|answer| A
%    B -.-> C["address<br/>203.0.113.7"]
%    A -->|then connect to the address| C
%```
%ENDMD
\begin{center}
\begin{tikzpicture}[>=Stealth]
  \node[nb] (a) at (0,0) {your code\\\code{getaddrinfo("api.x.test")}};
  \node[nb,draw=cpurple,fill=cpurple!10] (b) at (4.8,0) {DNS servers\\(distributed database)};
  \node[nb,draw=cgreen,fill=cgreen!10] (c) at (9.2,0) {address\\\code{203.0.113.7}};
  \draw[->] ([yshift=3pt]a.east) -- node[above,font=\tiny\sffamily]{name?} ([yshift=3pt]b.west);
  \draw[<-] ([yshift=-3pt]a.east) -- node[below,font=\tiny\sffamily]{answer} ([yshift=-3pt]b.west);
  \draw[->,dashed] (b) -- (c);
  \draw[->,thick,cblue] (c.south) -- ++(0,-0.7) -| node[pos=0.25,below,font=\tiny\sffamily]
    {then \code{connect((ip, 443))}} (a.south);
\end{tikzpicture}
\end{center}

The \code{socket} module has the functions. Names below resolve locally, so the output is
stable; the \code{getaddrinfo} tuple shape is the same for real hosts.

\begin{pycode}
import socket

print(socket.gethostbyname("localhost"))             # name -> IPv4 str
print(socket.getservbyname("http"), socket.getservbyport(443))
for info in socket.getaddrinfo("localhost", 80, socket.AF_INET,
                               socket.SOCK_STREAM):
    print(info)
\end{pycode}
\begin{rawcode}
127.0.0.1
80 https
(<AddressFamily.AF_INET: 2>, <SocketKind.SOCK_STREAM: 1>, 0, '', ('127.0.0.1', 80))
\end{rawcode}

\begin{center}
\begin{tabularx}{\linewidth}{@{}l X@{}}
\toprule
\textbf{Function} & \textbf{Returns} \\
\midrule
\code{gethostbyname(name)} & one IPv4 address as \code{str} \\
\code{gethostbyname\_ex(name)} & \code{(canonical name, [aliases], [addresses])} \\
\code{getaddrinfo(host, port, ...)} & list of tuples with everything needed to build a socket;
filter with family and socket type \\
\code{getservbyname("http")} / \code{getservbyport(80)} & service name to port number and back (IANA list) \\
\bottomrule
\end{tabularx}
\end{center}

\subsection{Python email modules}

\begin{center}
\begin{tabularx}{\linewidth}{@{}l l X@{}}
\toprule
\textbf{Module} & \textbf{Protocol} & \textbf{Use} \\
\midrule
\code{email} & -- & build and parse message objects \\
\code{smtplib} & SMTP & send mail \\
\code{poplib} & POP3 & download mail \\
\code{imaplib} & IMAP & read mail kept on the server \\
\bottomrule
\end{tabularx}
\end{center}

Building a message needs no server, so it can be run here:

\begin{pycode}
from email.message import EmailMessage

msg = EmailMessage()
msg["From"] = "bot@example.com"
msg["To"] = "me@example.com"
msg["Subject"] = "Nightly eval finished"
msg.set_content("Score: 0.91")
print(msg.as_string())
\end{pycode}
\begin{rawcode}
From: bot@example.com
To: me@example.com
Subject: Nightly eval finished
Content-Type: text/plain; charset="utf-8"
Content-Transfer-Encoding: 7bit
MIME-Version: 1.0

Score: 0.91
\end{rawcode}

Sending uses \code{smtplib.SMTP(host, 587)}, then \code{starttls()}, \code{login(...)},
\code{send\_message(msg)}. I did not run that: it needs a real mail server (syntax-checked
only). Never hard-code the password; read it from an environment variable.

\subsection{Other protocols}

\code{ftplib} (FTP) moves files; the standard library has modules for most other
internet protocols too. The point: before installing a package, check whether the
standard library already speaks the protocol.

% ------------------------------------------------------------
\section{Web Services and APIs}

Websites are made for human eyes. An \emph{API} serves the same data to programs in a
structured format, usually JSON. A scraper breaks whenever the page layout changes; an
API changes much less often. Most APIs need a registered \emph{key} (a long token) that
identifies you and lets the provider limit your traffic.

An LLM API call is exactly this: an HTTPS request with a key in a header, a JSON body,
and a JSON reply. The part beginners miss is the \emph{status code}.

%MD
%```mermaid
%sequenceDiagram
%    participant C as client
%    participant S as API server
%    C->>S: POST /v1/complete
%    S-->>C: 429 slow down
%    Note over C: wait, back off
%    C->>S: POST (retry)
%    S-->>C: 200 + JSON body
%```
%ENDMD
\begin{center}
\begin{tikzpicture}[>=Stealth]
  \node[font=\sffamily\bfseries\small,text=cblue] at (0,0.3) {client};
  \node[font=\sffamily\bfseries\small,text=cgreen] at (7,0.3) {API server};
  \draw[cgray] (0,0) -- (0,-4.6);
  \draw[cgray] (7,0) -- (7,-4.6);
  \draw[->,thick,cblue] (0,-0.6) -- node[above,font=\scriptsize\ttfamily]{POST /v1/complete} (7,-0.6);
  \draw[->,thick,cred]  (7,-1.4) -- node[above,font=\scriptsize\ttfamily]{429 slow down} (0,-1.4);
  \node[font=\scriptsize\sffamily,text=cgray,fill=white] at (0,-2.1) {wait, back off};
  \draw[->,thick,cblue] (0,-2.8) -- node[above,font=\scriptsize\ttfamily]{POST (retry)} (7,-2.8);
  \draw[->,thick,cgreen](7,-3.6) -- node[above,font=\scriptsize\ttfamily]{200 + JSON body} (0,-3.6);
\end{tikzpicture}
\end{center}

\begin{center}
\begin{tabularx}{\linewidth}{@{}l l X@{}}
\toprule
\textbf{Code} & \textbf{Meaning} & \textbf{What your code should do} \\
\midrule
200 & OK & parse the JSON body \\
400 & bad request & your payload is wrong; fix it, do not retry \\
401 / 403 & bad key / not allowed & stop; check the key \\
404 & not found & check the URL or model name \\
429 & too many requests & wait (see a \code{retry-after} header), then retry \\
500 / 503 & server error / overloaded & retry with exponential backoff \\
\bottomrule
\end{tabularx}
\end{center}

To keep this runnable offline, here is a tiny fake API on localhost using only the
standard library, and a client that returns the status code and parsed body:

\begin{pycode}
import json
import threading
import urllib.error
import urllib.request
from http.server import BaseHTTPRequestHandler, HTTPServer

class FakeLLM(BaseHTTPRequestHandler):
    def do_POST(self):
        size = int(self.headers.get("Content-Length", 0))
        body = json.loads(self.rfile.read(size))
        if self.headers.get("x-api-key") != "test-key":
            code, out = 401, {"error": "bad key"}
        elif len(body["prompt"]) > 20:
            code, out = 429, {"error": "slow down"}
        else:
            code, out = 200, {"text": body["prompt"][::-1]}
        data = json.dumps(out).encode("utf-8")
        self.send_response(code)
        self.send_header("Content-Type", "application/json")
        self.send_header("Content-Length", str(len(data)))
        self.end_headers()
        self.wfile.write(data)

    def log_message(self, *args):          # silence request logging
        pass

srv = HTTPServer(("127.0.0.1", 0), FakeLLM)
threading.Thread(target=srv.serve_forever, daemon=True).start()
url = f"http://127.0.0.1:{srv.server_address[1]}/v1/complete"

def call(prompt, key="test-key"):
    req = urllib.request.Request(
        url,
        data=json.dumps({"prompt": prompt}).encode("utf-8"),
        headers={"Content-Type": "application/json", "x-api-key": key},
        method="POST",
    )
    try:
        with urllib.request.urlopen(req) as resp:
            return resp.status, json.loads(resp.read())
    except urllib.error.HTTPError as err:   # 4xx and 5xx raise
        return err.code, json.loads(err.read())

print(call("hello"))
print(call("hello", key="wrong"))
print(call("a very long prompt that is over the limit"))
srv.shutdown()
\end{pycode}
\begin{rawcode}
(200, {'text': 'olleh'})
(401, {'error': 'bad key'})
(429, {'error': 'slow down'})
\end{rawcode}

\begin{typecheck}
The request body is \code{bytes} (\code{json.dumps(...).encode("utf-8")}). The reply
body is also \code{bytes}; \code{json.loads} accepts bytes and returns a \code{dict}. With
the popular \code{requests} package the same idea is shorter (\code{pip install requests};
syntax-checked only): \code{resp.status\_code} is an \code{int}, \code{resp.json()} is a
\code{dict}, and \code{resp.raise\_for\_status()} turns 4xx/5xx into an exception.
\end{typecheck}

\begin{pycode}
# needs: pip install requests
import os
import requests

resp = requests.post(
    "https://api.example.test/v1/complete",
    headers={"x-api-key": os.environ["MY_API_KEY"]},
    json={"prompt": "hello"},
    timeout=30,
)
if resp.status_code == 429:
    print("rate limited; retry after", resp.headers.get("retry-after"))
else:
    resp.raise_for_status()                    # 4xx/5xx -> HTTPError
    data = resp.json()                         # dict
\end{pycode}

\begin{gamerule}[: always set a timeout, keep keys out of code]
A network call without a timeout can hang forever. Put API keys in environment variables,
never in source files or git.
\end{gamerule}

% ------------------------------------------------------------
\section{Data Serialization}

Memory holds objects; a network or a file holds bytes. Converting one to the other is
\emph{serialization} (also \emph{marshaling}); the reverse is \emph{deserialization}.

%MD
%```mermaid
%flowchart LR
%    O1["Python object<br/>(dict)"] -->|dumps| B1["bytes"] --> W["wire / file"] --> B2["bytes"] -->|loads| O2["object"]
%```
%ENDMD
\begin{center}
\begin{tikzpicture}[>=Stealth]
  \node[nb,draw=cgreen,fill=cgreen!10] (o1) at (0,0) {Python object\\\code{dict}};
  \node[nb,draw=cgold,fill=cgold!15] (b1) at (4.2,0) {bytes};
  \node[nb,draw=cgray,fill=cgray!10] (w) at (7.4,0) {wire / file};
  \node[nb,draw=cgold,fill=cgold!15] (b2) at (10.2,0) {bytes};
  \node[nb,draw=cgreen,fill=cgreen!10] (o2) at (13.2,0) {object};
  \draw[->] (o1) -- node[above,font=\tiny\ttfamily]{dumps} (b1);
  \draw[->] (b1) -- (w);
  \draw[->] (w) -- (b2);
  \draw[->] (b2) -- node[above,font=\tiny\ttfamily]{loads} (o2);
\end{tikzpicture}
\end{center}

\subsection{JSON vs pickle}

\begin{pycode}
import datetime
import json
import pickle

record = {"model": "m1", "tokens": 120, "ok": True, "tags": ["a", "b"]}
wire = json.dumps(record).encode("utf-8")
print(type(wire), wire)
print(json.loads(wire.decode("utf-8")) == record)

stamp = datetime.datetime(2026, 1, 5, 9, 30)
try:
    json.dumps({"at": stamp})
except TypeError as err:
    print("TypeError:", err)
print(json.dumps({"at": stamp}, default=str))

blob = pickle.dumps(stamp)
print(type(blob), pickle.loads(blob) == stamp)

class Tiny:
    def __init__(self, n):
        self.n = n

copy = pickle.loads(pickle.dumps(Tiny(3)))
print(type(copy).__name__, copy.n)

print(json.loads(json.dumps((1, 2))))
try:
    json.dumps({1, 2})
except TypeError as err:
    print("TypeError:", err)
print(json.dumps({1: "a"}), json.loads(json.dumps({1: "a"})))
\end{pycode}
\begin{rawcode}
<class 'bytes'> b'{"model": "m1", "tokens": 120, "ok": true, "tags": ["a", "b"]}'
True
TypeError: Object of type datetime is not JSON serializable
{"at": "2026-01-05 09:30:00"}
<class 'bytes'> True
Tiny 3
[1, 2]
TypeError: Object of type set is not JSON serializable
{"1": "a"} {'1': 'a'}
\end{rawcode}

\begin{typecheck}
JSON only knows: \code{dict} (string keys), \code{list}, \code{str}, \code{int}/\code{float},
\code{True}/\code{False}, \code{None}. Everything else fails or changes shape:
a \code{tuple} comes back as a \code{list}, a \code{set} or \code{datetime} raises
\code{TypeError}, an \code{int} dict key comes back as \code{str}. Fix unknown types with
\code{default=str} or by converting yourself. \code{pickle.dumps} returns \code{bytes},
\code{json.dumps} returns \code{str}.
\end{typecheck}

\begin{trap}[: never unpickle data you do not trust]
\code{pickle.loads} can run arbitrary code. Data from a network, an upload or a shared
cache must never go through pickle. Use JSON (or another data-only format) at trust
boundaries. Pickle is fine for your own temporary files; \code{multiprocessing} uses it
internally.
\end{trap}

\subsection{Other formats}

\begin{center}
\begin{tabularx}{\linewidth}{@{}l l X@{}}
\toprule
\textbf{Format} & \textbf{Text or binary} & \textbf{Notes} \\
\midrule
JSON & text & human-readable, universal on the web, few types \\
pickle & binary & any Python object; Python-only; unsafe for untrusted data \\
MessagePack, msgspec & binary & JSON-like data, smaller and faster \\
Protocol Buffers, Avro, Thrift & binary & schema-defined, used with RPC \\
\bottomrule
\end{tabularx}
\end{center}

Packages such as \textbf{Pydantic} and marshmallow go the other way: they turn raw data
into typed objects and \emph{validate} it (types, ranges, required fields). Web servers
use them so bytes from the wire land in the right structures. Pydantic is the one you
will meet most in AI work (structured LLM output).

% ------------------------------------------------------------
\section{Remote Procedure Calls}

An RPC makes a function on another machine look like a local call. Your side serializes
the arguments, sends bytes, and the remote side decodes, runs the function, and sends the
encoded result back.

%MD
%```mermaid
%flowchart LR
%    A["proxy.double(7)"] --> B["encode arguments"] --> C["network"] --> D["decode, find function"] --> E["run double(7)"]
%    E -.->|encode result, send back| A
%```
%ENDMD
\begin{center}
\begin{tikzpicture}[>=Stealth]
  \node[nb] (a) at (0,0) {\code{proxy.double(7)}};
  \node[nb,draw=cgold,fill=cgold!15] (b) at (4.2,0) {encode\\arguments};
  \node[nb,draw=cgray,fill=cgray!10] (c) at (8,0) {network};
  \node[nb,draw=cgold,fill=cgold!15] (d) at (11.8,0) {decode,\\find function};
  \node[nb,draw=cgreen,fill=cgreen!10] (e) at (11.8,-1.5) {run \code{double(7)}};
  \draw[->] (a) -- (b);
  \draw[->] (b) -- (c);
  \draw[->] (c) -- (d);
  \draw[->] (d) -- (e);
  \draw[->,dashed,cblue] (e.west) -- ++(-2.4,0) -- node[below,font=\tiny\sffamily]
    {encode result, send back} (a.south |- e.west) -- (a.south);
\end{tikzpicture}
\end{center}

\subsection{XML-RPC}

The standard library has one RPC: XML-RPC over HTTP. This runs fully on localhost:

\begin{pycode}
import threading
import xmlrpc.client
from xmlrpc.server import SimpleXMLRPCServer

def cost_usd(tokens_in, tokens_out):
    return tokens_in * 1e-7 + tokens_out * 5e-7

server = SimpleXMLRPCServer(("127.0.0.1", 0), logRequests=False)
port = server.server_address[1]
server.register_function(cost_usd, "cost_usd")
t = threading.Thread(target=server.serve_forever, daemon=True)
t.start()

proxy = xmlrpc.client.ServerProxy(f"http://127.0.0.1:{port}/")
result = proxy.cost_usd(2000, 500)       # looks like a local call
print(result, type(result))
try:
    proxy.nope()                         # no such function on the server
except xmlrpc.client.Fault as err:
    print("Fault:", err.faultCode)
server.shutdown()
\end{pycode}
\begin{rawcode}
0.00045 <class 'float'>
Fault: 1
\end{rawcode}

\code{proxy.cost\_usd} does not exist in the client; the proxy turns any attribute
access into a remote call. Remote errors arrive as \code{xmlrpc.client.Fault}.

\subsection{Other RPC systems}

\begin{center}
\begin{tabularx}{\linewidth}{@{}l l X@{}}
\toprule
\textbf{System} & \textbf{Package} & \textbf{Notes} \\
\midrule
XML-RPC & stdlib \code{xmlrpc} & XML over HTTP, simple, verbose \\
JSON-RPC & \code{jsonrpcserver}, \code{jsonrpcclient} & same idea in JSON \\
MessagePack RPC & \code{msgpack-rpc-python} & binary, uses Tornado \\
zerorpc & \code{zerorpc} & ZeroMQ + MessagePack; methods of a class become endpoints \\
gRPC & \code{grpcio}, \code{grpcio-tools} & Google; Protocol Buffers; streaming; fast \\
\bottomrule
\end{tabularx}
\end{center}

gRPC works from a schema. You describe the service in a \code{.proto} file, and a tool
generates Python client and server stubs:

\begin{rawcode}
service Pricing {
  rpc Cost (CostRequest) returns (CostReply);
}
message CostRequest { int32 tokens_in = 1; int32 tokens_out = 2; }
message CostReply   { double usd = 1; }
\end{rawcode}

A rule of thumb: REST/JSON for public APIs, gRPC for fast service-to-service calls
inside your own system.

% ------------------------------------------------------------
\section{Remote Management}

Tools that run commands on many machines at once.

\begin{center}
\begin{tabularx}{\linewidth}{@{}l X@{}}
\toprule
\textbf{Tool} & \textbf{Notes} \\
\midrule
Salt & written in Python; ZeroMQ based; scales to thousands of machines \\
Ansible & written in Python; agentless, uses SSH \\
Puppet, Chef & popular, tied to Ruby \\
\bottomrule
\end{tabularx}
\end{center}

% ------------------------------------------------------------
\section{Big Fat Data}

One machine, even a big one, hits limits in disk, memory and CPU. Spreading data across
many machines can be faster than a single fast one. The core idea is \emph{MapReduce}:
split the data, process each piece independently (map), then combine the partial results
(reduce).

%MD
%```mermaid
%flowchart LR
%    I["big input"] --> M1["map: part 1"] --> P1["partial 1"] --> R["reduce"]
%    I --> M2["map: part 2"] --> P2["partial 2"] --> R
%    I --> M3["map: part 3"] --> P3["partial 3"] --> R
%    R --> RES["result"]
%```
%ENDMD
\begin{center}
\begin{tikzpicture}[n/.style={nb,minimum width=1.7cm,font=\scriptsize\sffamily},>=Stealth]
  \node[n,draw=cgray,fill=cgray!10] (in) at (0,0) {big input};
  \foreach \y/\i in {1/1,0/2,-1/3} {
    \node[n,draw=cblue,fill=cblue!10] (m\i) at (3.4,\y) {map: part \i};
    \draw[->] (in) -- (m\i);
    \node[n,draw=cgold,fill=cgold!15] (p\i) at (6.4,\y) {partial \i};
    \draw[->] (m\i) -- (p\i);
  }
  \node[n,draw=cgreen,fill=cgreen!10] (r) at (9.4,0) {reduce};
  \foreach \i in {1,2,3} \draw[->] (p\i) -- (r);
  \node[n,draw=cgreen,fill=cgreen!10,right=8mm of r] (res) {result};
  \draw[->] (r) -- (res);
\end{tikzpicture}
\end{center}

\begin{center}
\begin{tabularx}{\linewidth}{@{}l X@{}}
\toprule
\textbf{Tool} & \textbf{Notes} \\
\midrule
Hadoop & Java; copies data between machines and writes to disk at each step, so batch jobs are slow. Hadoop streaming lets any language, including Python, plug in \\
Spark & much faster (in memory); Python API (PySpark); the usual choice today \\
Disco & Python + Erlang MapReduce; little maintained \\
Dask & written in Python; works with NumPy, pandas, scikit-learn; scales from a laptop to clusters \\
\bottomrule
\end{tabularx}
\end{center}

``Big data'' often just means ``too big for my machine''. Try a bigger machine or a
smarter format first; use a cluster only when you must.

% ------------------------------------------------------------
\section{Clouds}

Servers fail, and with thousands of them, failure is routine. Peter Deutsch's ``fallacies
of distributed computing'' list the false assumptions that cause bugs.

\begin{center}
\begin{tabularx}{\linewidth}{@{}X X@{}}
\toprule
\multicolumn{2}{@{}l}{\textbf{Things that are NOT true on a network}} \\
\midrule
The network is reliable. & Latency is zero. \\
Bandwidth is infinite. & The network is secure. \\
Topology does not change. & There is one administrator. \\
Transport cost is zero. & The network is homogeneous. \\
\bottomrule
\end{tabularx}
\end{center}

This is why calls to remote APIs need timeouts, retries and error handling (the status
code table above).

At scale, servers are treated as interchangeable (``livestock''): replace broken ones
rather than nursing them. Renting machines in the cloud turns that into an API call.

\begin{center}
\begin{tabularx}{\linewidth}{@{}l l X@{}}
\toprule
\textbf{Service} & \textbf{Python tool} & \textbf{Notes} \\
\midrule
Amazon Web Services & \code{boto3} & the market leader; replaces the web console with code \\
Google Cloud & Google client libraries & Google uses Python heavily \\
Microsoft Azure & Azure SDK for Python & third big vendor \\
OpenStack & its REST APIs & open source cloud framework \\
\bottomrule
\end{tabularx}
\end{center}

\subsection{Docker and Kubernetes}

A \emph{container} packages your app with its libraries, sharing the host's operating
system kernel. It is lighter than a virtual machine and heavier than a virtualenv.
\emph{Kubernetes} runs many containers for you: restarts failures, balances load, scales
up and down.

%MD
%```text
%Virtual machines                    Containers
%+----------+ +----------+           +----------+ +----------+
%| your app | | your app |           | your app | | your app |
%| libraries| | libraries|           | libraries| | libraries|
%| guest OS | | guest OS |           +----------+-+----------+
%+----------+-+----------+           | container runtime     |
%| hypervisor             |           +-----------------------+
%+------------------------+           | host OS (shared kernel)|
%| host OS and hardware   |           +-----------------------+
%+------------------------+
%```
%ENDMD
\begin{center}
\begin{tikzpicture}[l/.style={draw,minimum width=3.3cm,minimum height=6.5mm,font=\scriptsize\sffamily,
    align=center,inner sep=2pt}]
  \node[font=\sffamily\bfseries\small] at (0,3.4) {Virtual machines};
  \foreach \x in {-1.8,1.8} {
    \node[l,fill=cgreen!12,draw=cgreen] at (\x,2.6) {your app};
    \node[l,fill=cgold!15,draw=cgold]   at (\x,2.0) {libraries};
    \node[l,fill=cred!10,draw=cred]     at (\x,1.4) {guest OS (full)};
  }
  \node[l,minimum width=7.2cm,fill=cpurple!12,draw=cpurple] at (0,0.8) {hypervisor};
  \node[l,minimum width=7.2cm,fill=cgray!12,draw=cgray] at (0,0.2) {host OS and hardware};
  \begin{scope}[xshift=8.4cm]
  \node[font=\sffamily\bfseries\small] at (0,3.4) {Containers};
  \foreach \x in {-1.8,1.8} {
    \node[l,fill=cgreen!12,draw=cgreen] at (\x,2.0) {your app};
    \node[l,fill=cgold!15,draw=cgold]   at (\x,1.4) {libraries};
  }
  \node[l,minimum width=7.2cm,fill=cpurple!12,draw=cpurple] at (0,0.8) {container runtime};
  \node[l,minimum width=7.2cm,fill=cgray!12,draw=cgray] at (0,0.2) {host OS (shared kernel)};
  \end{scope}
\end{tikzpicture}
\end{center}

Python clients: \code{docker} and \code{kubernetes} (\code{pip install docker kubernetes}).
Docker is how most AI services get shipped.

% ------------------------------------------------------------
\section{Practice}

\begin{practice}
\begin{enumerate}[leftmargin=*,itemsep=2pt]
\item Build a TCP time service on localhost: when a client sends \code{time}, the server
  replies with the current time as an ISO string. Use the length-prefix helpers and port 0.
\item Extend the fake API above: make it return 503 for the first two calls, then 200.
  Write a client that retries with doubling waits and gives up after four tries.
\item Add prefix matching to the in-process \code{Broker}, so a subscriber to
  \code{"eval."} also receives \code{"eval.done"}.
\item Write an XML-RPC server with one function that returns a dict. Call it from a client.
  What happens if the function returns \code{None}?
\item Explain in two sentences why \code{pickle.loads} on bytes from a socket is dangerous,
  and which format you would use instead.
\item (Needs Redis or Valkey.) Push ten job strings to a list, then write a worker that
  loops on \code{brpop}, sleeps 0.5 s per job, and prints how many jobs remain.
\end{enumerate}
\end{practice}

\begin{achievement}{Networking Basics}
You can now:
\begin{itemize}[leftmargin=*,itemsep=1pt]
\item explain the difference between UDP and TCP and write a localhost echo with sockets;
\item frame messages on a TCP stream and keep \code{bytes} and \code{str} apart;
\item pick request-reply, push or pub-sub for a given problem;
\item read HTTP status codes from an API and decide whether to retry;
\item choose between JSON and pickle, and explain the pickle danger.
\end{itemize}
\end{achievement}
